\documentclass[12pt]{article}
\usepackage[margin=1.05in]{geometry}
\usepackage{amsmath,amssymb,amsthm,mathtools}
\usepackage{booktabs,array,longtable,microtype}
\usepackage{enumitem}
\usepackage{hyperref}
\hypersetup{colorlinks=true,linkcolor=blue,citecolor=blue,urlcolor=blue}

\newtheorem{theorem}{Theorem}[section]
\newtheorem{proposition}[theorem]{Proposition}
\newtheorem{corollary}[theorem]{Corollary}
\newtheorem{definition}[theorem]{Definition}
\newtheorem*{remark}{Remark}

\newcommand{\F}{\mathbb{F}}
\newcommand{\Z}{\mathbb{Z}}
\newcommand{\Weyl}{\mathrm{W}}
\newcommand{\Spec}{\mathrm{Spec}}
\newcommand{\rankop}{\mathrm{rank}}
\newcommand{\Xmin}{X_{\min}}
\newcommand{\Zmin}{Z_{\min}}
\newcommand{\Roots}{\Phi}

\title{The W33 Holographic Code Tower and the Minimal Logical $E_6$ Surface\\[4pt]
\large Algebraic-Geometry Codes, Qutrit CSS Homology, and a Signed Rank-$81$ Phase Frame}
\author{Wil Dahn\\[3pt]\small Independent Research\\[-1pt]\small \texttt{github.com/wilcompute/W33-Theory}}
\date{Revised May 22, 2026}

\begin{document}
\maketitle

\begin{abstract}
This paper gives a corrected and strengthened formulation of the W33 holographic code tower.
The construction has two linked substrates: the internal generalized-quadrangle graph
$W(3,3)$, with $40$ vertices and $240$ edges, and the boundary horizon graph $K_{12}$,
whose orientable embedding genus is
\[
        g(K_{12})=\frac{(12-3)(12-4)}{12}=6.
\]
The first layer is a tower of ternary evaluation/CSS codes organized by the triple
$(q,g,h)=(3,6,12)$, including the bulk code $[[240,81,3]]_3$ and the boundary code
$[72,66,3]_3$.  The second layer is the minimal logical surface of the canonical
W33 edge CSS code.  Its projective minimal logical rays satisfy
\[
|\Xmin|=160,\qquad |\Zmin|=1620,
\]
with support biregularity $160\cdot81=1620\cdot8=12960$.  New support-geometry analysis shows that these supports are not opaque: the $160$ minimal $X$ supports are isotropic line-stars and the $1620$ minimal $Z$ supports are ordinary quadrangles of the $W(3,3)$ collinearity graph.  At the vector level,
\[
|\Xmin^{\F_3}|=320,\qquad |\Zmin^{\F_3}|=3240,
\]
and the nonzero symplectic-pairing census is exactly
\[
        25920+25920=51840=|\Weyl(E_6)|.
\]
Retaining the actual $\F_3$ phase gives a signed matrix
$A\in\{-1,0,1\}^{160\times1620}$ satisfying
\[
        AA^T = 160\Pi,
        \qquad \rank(\Pi)=81,
        \qquad \Pi^2=\Pi.
\]
Thus the unsigned shadow of the minimal logical surface is the Weyl order of $E_6$,
while the signed phase geometry is an exact rank-$81$ projector, matching the protected
homology dimension of the W33 bulk code.  We further record the $3$-adic $X$-visibility
scheme with overlaps $1,3,9,27$, the bounded dual $Z$-visibility scheme with overlaps
$0,1,2,3,4$, and the toroidal metric generating function
\[
P(t)=68+147t+127t^2+86t^3+54t^4+19t^5+3t^6=(1+t)Q(t),
\]
where $Q(t)=68+79t+48t^2+38t^3+16t^4+3t^5$ encodes the phase-kernel number $79$,
the genus numerator $12$, the middle eigenvalue $72$, and the Eisenstein norm
$11^2\cdot21$.  The result is a single finite qutrit object in which $E_6$ appears
not merely as a dimension match but as the exact noncommutation census of minimal
logical errors.
\end{abstract}

\tableofcontents

\bigskip
\noindent\textbf{MSC 2020.}
94B27, 05E30, 17B25, 81P70, 81T30.

\medskip
\noindent\textbf{Keywords.}
$W(3,3)$, generalized quadrangle, qutrit CSS code, minimal logical errors, $E_6$,
phase frame, algebraic geometry codes, holographic codes, Schlaefli geometry.

\newpage

% ===========================================================
\section{Executive summary}
% ===========================================================

The central claim is finite and testable.  The internal substrate is not merely a list of
coincidences; it is the symplectic generalized quadrangle $W(3,3)$ with collinearity graph
\[
        \operatorname{srg}(40,12,2,4).
\]
The $240$ graph edges support a qutrit CSS code whose protected sector has dimension $81$.
The new theorem stack says that the minimal logical error surface simultaneously sees
\[
        |\Weyl(E_6)|=51840
\]
when signs are forgotten and
\[
        \dim H_1=81
\]
when the $\F_3$ phase is retained.  This is the sharpest bridge in the present paper:
\[
\boxed{\text{forget phase}\Rightarrow E_6\text{ Weyl noncommutation count}},
\]
\[
\boxed{\text{retain phase}\Rightarrow \text{rank-}81\text{ homological projector}}.
\]

% ===========================================================
\section{Two substrates, not one}
% ===========================================================

A major correction is to keep two finite objects distinct.

\begin{definition}[Internal W33 substrate]
The internal substrate is the collinearity graph of the symplectic generalized quadrangle
$W(3,3)$.  Its vertices are the $40$ projective points of $\F_3^4$, and two vertices are
adjacent exactly when their symplectic form vanishes.  The resulting graph is
\[
        \operatorname{srg}(40,12,2,4),
\]
so it has $40$ vertices, valency $12$, $240$ edges, $\lambda=2$, and $\mu=4$.
\end{definition}

\begin{definition}[Horizon substrate]
The horizon substrate is the complete graph $K_{12}$.  It has
\[
        |E(K_{12})|=\binom{12}{2}=66
\]
and orientable complete-graph genus
\[
        \gamma(K_{12})=\frac{(12-3)(12-4)}{12}=6.
\]
The triple
\[
        (q,g,h)=(3,6,12)
\]
will always mean: ternary field, horizon genus, and horizon valency/Coxeter number.
\end{definition}

\begin{remark}
The internal object supplies the $240$ physical edge-qutrits and the $81$ protected
homology sector.  The horizon object supplies the boundary count $66=\binom{12}{2}$
and the genus $6$.  Confusing these two objects makes the numerology look stronger but
the mathematics weaker; separating them makes the theorem stack cleaner.
\end{remark}

% ===========================================================
\section{The corrected code tower}
% ===========================================================

The tower is organized by the triple $(q,g,h)=(3,6,12)$.  The bulk and horizon lengths are
\[
        n_B=q^5-q=3^5-3=240,
        \qquad
        n_H=gh=72.
\]
The bulk length factors as
\[
        240=3(3^4-1)=3(3-1)(3+1)(3^2+1)=3\cdot2\cdot4\cdot10.
\]

\begin{center}
\renewcommand{\arraystretch}{1.2}
\begin{tabular}{cllrrr}
\toprule
Layer & Interpretation & Code/data & $n$ & $k$ & identity\\
\midrule
$6$ & affine completion & $[243,237,3]_3$ & $243$ & $237$ & $3^5$\\
$5$ & W33 edge bulk & $[[240,81,3]]_3$ & $240$ & $81$ & $3^5-3$\\
$4$ & $E_7/E_6$ gap & $[55,49,3]_3$ & $55$ & $49$ & $133-78$\\
$3$ & $E_7/(E_6\times U(1))$ & $[54,48,3]_3$ & $54$ & $48$ & $133-79$\\
$2$ & Spin$(10)$ spinor & $[32,26,3]_3$ & $32$ & $26$ & $2^5$\\
$1$ & $K_{12}$ horizon & $[72,66,3]_3$ & $72$ & $66$ & $6\cdot11$\\
$0$ & entanglement wedge & kernel & -- & $15$ & $81-66$\\
\bottomrule
\end{tabular}
\end{center}

The boundary and bulk rates are
\[
        R_H={66\over72}={11\over12},
        \qquad
        R_B={81\over240},
\]
so
\[
        {R_H\over R_B}=\frac{11/12}{81/240}=\frac{220}{81}.
\]
The numerator $220=\binom{12}{3}$ is the number of triangles in the horizon graph.

% ===========================================================
\section{Riemann--Roch dimension layer}
% ===========================================================

\begin{proposition}[Dimension rule]
For a genus-$6$ evaluation code $C_{\mathcal L}(D,G)$ with $\deg(G)=n-1>2g-2$ and
$\ell(K-G)=0$,
\[
        k=\ell(G)=\deg(G)-g+1=n-g.
\]
Hence $n-k=g=6$.
\end{proposition}

\begin{proof}
Riemann--Roch gives
\[
        \ell(G)-\ell(K-G)=\deg(G)-g+1.
\]
When $\deg(G)>2g-2$, the correction term vanishes.  With $g=6$ and $\deg(G)=n-1$,
this gives $k=n-6$.
\end{proof}

\begin{remark}[Distance honesty boundary]
The dimension rule explains the repeated value $n-k=6$ in the evaluation layers.  It does
not by itself prove distance $3$.  Distance statements are finite certificates or direct
constructions, not automatic consequences of the displayed Riemann--Roch calculation.
\end{remark}

% ===========================================================
\section{Exceptional ladder identities}
% ===========================================================

\begin{proposition}[Logical ladder]
The logical counts satisfy
\begin{align*}
15&=\dim(G_2)+1,\\
26&=\dim(E_6)-\dim(F_4),\\
48&=|\Roots(F_4)|,\\
49&=\operatorname{rank}(E_7)^2,\\
66&=\binom{12}{2}=g(h-1),\\
81&=3^4=q^{\operatorname{rank}(F_4)},\\
237&=3(78+1)=q\dim(E_6\times U(1)).
\end{align*}
Moreover,
\begin{align*}
81-48 &= 48-15 = 33=q(h-1),\\
49-48 &= 1,\\
81-49 &= 32=2^5,\\
26-15 &= 11=h-1,\\
(81-66)+(66-49)+(49-48)+(48-26)+(26-15)&=66.
\end{align*}
\end{proposition}

\begin{proof}
All identities follow from the integer values in Proposition~\ref{prop:ladder}.  The point
of the proposition is not the arithmetic alone, but that the same small set of operators
$q=3$, $g=6$, $h=12$, and the exceptional chain $G_2,F_4,E_6,E_7,E_8$ generates the entire
ladder.
\end{proof}

% ===========================================================
\section{The canonical W33 edge CSS surface}
% ===========================================================

We now move from the evaluation-code tower to the canonical edge CSS surface.  Let
$C_1\cong\F_3^{240}$ be the ternary vector space on the W33 edges.  The code is a CSS
code on these $240$ qutrits with protected homology dimension
\[
        k_B=81.
\]
The minimal logical surface is the set of minimal nontrivial $X$ and $Z$ representatives,
modulo scalar multiplication by $\F_3^\times=\{1,2\}$ unless explicitly stated otherwise.

\begin{definition}[Projective and vector minimal logical sets]
Let $\Xmin$ and $\Zmin$ be the projective minimal logical rays, and let
$\Xmin^{\F_3}$ and $\Zmin^{\F_3}$ denote the corresponding nonzero vectors after retaining
the two scalar representatives in each projective ray.
\end{definition}

\begin{theorem}[Minimal logical census]
\label{thm:mincensus}
The canonical W33 edge CSS code has minimal distances
\[
        d_X=3,\qquad d_Z=4,
\]
and minimal logical counts
\[
        |\Xmin|=160,
        \qquad
        |\Xmin^{\F_3}|=320,
\]
\[
        |\Zmin|=1620,
        \qquad
        |\Zmin^{\F_3}|=3240.
\]
\end{theorem}

\begin{proof}[Certificate proof]
The theorem is a finite enumeration certificate: the repository enumeration constructs all
minimal nontrivial logical representatives and removes scalar duplicates.  The scalar
multiplication check gives the factors $320=2\cdot160$ and $3240=2\cdot1620$.  The
minimal weights are the least nontrivial weights observed in the certified $X$ and $Z$
logical classes.
\end{proof}

\begin{remark}[Correcting the old $6480$ count]
The number $6480$ should not be read as the number of unique $\F_3$ vectors.  It is an
oriented-walk presentation count.  The unique vector count after quotienting presentations
is $3240$ for minimal $Z$ vectors.
\end{remark}

\begin{theorem}[Minimal support geometry theorem]
\label{thm:supportgeom}
The projective supports in Theorem~\ref{thm:mincensus} have the following intrinsic
finite-geometric model.
\begin{enumerate}[label=(\roman*)]
\item Each minimal $X$ support is an \emph{isotropic line-star}: choose one of the
$40$ totally isotropic lines $L\cong K_4$ of $W(3,3)$ and one point $p\in L$; the
support is the three graph-edges joining $p$ to the other points of $L$.  Hence
\[
        |\Xmin|=40\cdot4=160.
\]
\item Each minimal $Z$ support is an \emph{ordinary quadrangle}: choose a noncollinear
point pair $\{a,b\}$ and two of its four common neighbours $c,d$, producing the
4-cycle $a-c-b-d-a$.  Since the collinearity graph has $\binom{40}{2}-240=540$
nonedges and each quadrangle is counted by its two diagonals,
\[
        |\Zmin|={540\binom42\over2}=1620.
\]
\end{enumerate}
With the natural $\F_3$ cycle coefficients, this support model reproduces the
$81/8$ biregularity, the $1,3,9,27$ $X$-overlap scheme, and the signed rank-$81$
projector.
\end{theorem}

\begin{proof}[Enumeration certificate]
The certificate constructs $W(3,3)$ from projective points of $\F_3^4$ with the
standard alternating form, builds its $240$ collinearity edges and $40$ isotropic
$K_4$ lines, and then enumerates all line-supported weight-$3$ $X$ vectors satisfying
$d_2^T x=0$.  Exactly the $160$ line-stars survive.  It also enumerates all ordinary
quadrangles by the SRG parameter $\mu=4$ and orients each 4-cycle over $\F_3$.
The resulting signed projective pairing matrix has spectrum $160^{81}\oplus0^{79}$,
and its unsigned incidence matrix has off-diagonal $X$-overlap distribution
$1^{6480},3^{4320},9^{1440},27^{480}$.  Thus the support model is not merely
cardinality-matched; it reproduces the full minimal-logical scheme.
\end{proof}

% ===========================================================
\section{The $E_6$ noncommutation theorem}
% ===========================================================

Let $\langle x,z\rangle\in\F_3$ be the CSS symplectic pairing between an $X$ representative
and a $Z$ representative.  In projective language, $x$ and $z$ are visible to each other if
\[
        \langle x,z\rangle\ne0.
\]

\begin{theorem}[Support biregularity]
\label{thm:bireg}
The projective support incidence between minimal $X$ and $Z$ rays is biregular:
\[
        \deg_X=81,
        \qquad
        \deg_Z=8.
\]
Equivalently,
\[
        160\cdot81=1620\cdot8=12960.
\]
\end{theorem}

\begin{proof}
Counting visible projective pairs from the $X$ side gives $160\cdot81$.  Counting the same
incidence set from the $Z$ side gives $1620\cdot8$.  Both equal $12960$.
\end{proof}

\begin{theorem}[Minimal logical $E_6$ noncommutation theorem]
\label{thm:e6noncomm}
At the vector level, the $320\times3240$ ordered minimal $X/Z$ pairs have phase counts
\[
        0:984960,
        \qquad
        1:25920,
        \qquad
        2:25920.
\]
Therefore
\[
        \#\{(x,z):\langle x,z\rangle\ne0\}=25920+25920=51840=|\Weyl(E_6)|.
\]
Equivalently,
\[
        320\cdot162=3240\cdot16=51840.
\]
\end{theorem}

\begin{proof}
The phase count is obtained by retaining both nonzero scalar multiples of each projective
minimal ray and evaluating the $\F_3$ symplectic pairing on all ordered pairs.  The two
nonzero phases occur equally often.  Their sum is $51840$, the Weyl-group order of $E_6$.
Thus $E_6$ appears as the exact number of elementary noncommuting minimal logical
interactions.
\end{proof}

% ===========================================================
\section{The signed phase-frame theorem}
% ===========================================================

Define the signed projective phase matrix
\[
A_{xz}=\begin{cases}
0,& \langle x,z\rangle=0,\\
+1,& \langle x,z\rangle=1,\\
-1,& \langle x,z\rangle=2=-1.
\end{cases}
\]
Then $A\in\{-1,0,+1\}^{160\times1620}$.

\begin{theorem}[Minimal logical phase-frame theorem]
\label{thm:phaseframe}
The signed projective phase matrix satisfies
\[
        \Spec(AA^T)=160^{81}\oplus0^{79}.
\]
Equivalently,
\[
        (AA^T)^2=160AA^T.
\]
Thus
\[
        \Pi={1\over160}AA^T
\]
is an exact rank-$81$ projector.
\end{theorem}

\begin{proof}
The direct matrix certificate computes $S=AA^T$ over the integers and checks the exact
identity $S^2=160S$.  Therefore the only possible eigenvalues are $0$ and $160$.  The trace
is
\[
        \operatorname{tr}(S)=160\cdot81,
\]
so the multiplicity of the eigenvalue $160$ is $81$, and the remaining $79$ eigenvalues
are zero.
\end{proof}

\begin{corollary}[Unsigned and signed shadows]
The same minimal logical surface carries two complementary invariants:
\[
        \text{forget phase}\Rightarrow |\Weyl(E_6)|=51840,
\]
\[
        \text{retain phase}\Rightarrow \rank(A)=81.
\]
\end{corollary}

% ===========================================================
\section{Visibility schemes}
% ===========================================================

Let $U$ be the unsigned projective incidence matrix:
\[
        U_{xz}=1 \Longleftrightarrow \langle x,z\rangle\ne0.
\]

\begin{theorem}[$3$-adic $X$-visibility scheme]
\label{thm:xscheme}
Each minimal $X$ ray sees exactly $81$ minimal $Z$ rays.  For two distinct minimal $X$
rays, the number of shared visible $Z$ rays is always one of
\[
        1,3,9,27=3^0,3^1,3^2,3^3.
\]
For each fixed $X$ ray, the off-diagonal multiplicities are
\[
        1^{81},\qquad 3^{54},\qquad 9^{18},\qquad 27^6.
\]
Equivalently, the global unordered off-diagonal pair counts are
\[
        1:6480,
        \qquad
        3:4320,
        \qquad
        9:1440,
        \qquad
        27:480.
\]
The spectrum of $UU^T$ is
\[
648^1,
\quad
(144+36\sqrt6)^{24},
\quad
72^{30},
\quad
(144-36\sqrt6)^{24},
\quad
40^{81}.
\]
\end{theorem}

\begin{proof}
This is the exact Gram-certificate of the matrix $U$.  The diagonal entries are row degrees
$81$.  The off-diagonal entries of $UU^T$ have precisely the four listed values and
multiplicities.  Direct characteristic-polynomial/eigenvalue computation gives the displayed
spectrum.
\end{proof}

\begin{theorem}[Dual $Z$-visibility scheme]
\label{thm:zscheme}
Each minimal $Z$ ray sees exactly $8$ minimal $X$ rays.  For each fixed $Z$ ray, its overlap
distribution with all other $Z$ rays is
\[
        0^{1187},\qquad 1^{288},\qquad 2^{96},\qquad 3^{32},\qquad 4^{16}.
\]
Therefore each $Z$ ray has
\[
        288+96+32+16=432=16\cdot27
\]
nonzero $Z$-neighbours.
\end{theorem}

\begin{proof}
This is the dual Gram-certificate of $U^TU$.  The diagonal is $8$, and the off-diagonal
distribution is exactly the displayed distribution.  Summing the nonzero multiplicities gives
$432$.
\end{proof}

% ===========================================================
\section{Toroidal metric generating function}
% ===========================================================

The toroidal edge-multiplicity histogram
\[
        1^{12},2^{48},4^4,5^1,6^3
\]
has binomial moments
\[
        B_0,\\dots,B_6=68,147,127,86,54,19,3.
\]
Package these as
\[
        P(t)=68+147t+127t^2+86t^3+54t^4+19t^5+3t^6.
\]

\begin{theorem}[Toroidal metric generating function theorem]
\label{thm:toroidal}
The polynomial factors as
\[
        P(t)=(1+t)Q(t),
\]
where
\[
        Q(t)=68+79t+48t^2+38t^3+16t^4+3t^5.
\]
It satisfies
\[
        P(-1)=0,
        \qquad
        Q(-1)=12,
\]
\[
        P(1)=504=7\cdot72,
        \qquad
        Q(1)=252=21\cdot12.
\]
Modulo $\Phi_3(t)=t^2+t+1$,
\[
        P(t)\equiv11+55t,
\]
and therefore
\[
        N_{\Z[\omega]}(11+55\omega)=11^2\cdot21,
        \qquad \omega^2+\omega+1=0.
\]
\end{theorem}

\begin{proof}
Polynomial division gives $P(t)=(1+t)Q(t)$.  The evaluations are direct substitution.
Modulo $t^2+t+1$, reduce powers of $t$ by $t^2=-t-1$, obtaining $11+55t$.  The Eisenstein
norm is
\[
        N(a+b\omega)=a^2-ab+b^2,
\]
so
\[
        N(11+55\omega)=11^2-11\cdot55+55^2=2541=11^2\cdot21.
\]
\end{proof}

% ===========================================================
\section{Schlaefli and Standard Model bridge}
% ===========================================================

The W33 tower aligns with the standard exceptional breaking chain
\[
E_8\supset E_6\times SU(3)\supset \mathrm{Spin}(10)\supset SU(5)
\supset SU(3)\times SU(2)\times U(1).
\]
The $600$-cell $\{3,3,5\}$ has
\[
        V=120,
        \qquad
        E=720,
        \qquad
        F=1200,
        \qquad
        C=600.
\]
The clean bridge is
\[
        E(\{3,3,5\})=720=3\cdot240=q n_B.
\]
Also
\[
        F(\{3,3,5\})=1200=5n_B,
        \qquad
        V(\{3,3,5\})=120=12\cdot10.
\]

% ===========================================================
\section{Synthesis}
% ===========================================================

The current finite theorem stack is
\[
\boxed{d_X=3,\quad d_Z=4,}
\]
\[
\boxed{|\Xmin|=160,\quad |\Zmin|=1620,}
\]
\[
\boxed{160\cdot81=1620\cdot8=12960,}
\]
\[
\boxed{\#\{(x,z):\langle x,z\rangle\ne0\}=51840=|\Weyl(E_6)|,}
\]
\[
\boxed{AA^T/160\text{ is an exact rank-}81\text{ projector},}
\]
\[
\boxed{X\text{-overlaps}=1,3,9,27,\qquad Z\text{-overlaps}=0,1,2,3,4,}
\]
\[
\boxed{P(t)=(1+t)Q(t),\quad [t]Q(t)=79,\quad Q(-1)=12,\quad P(1)/7=72.}
\]
The strongest interpretation is: the canonical W33 edge CSS surface is a finite qutrit
substrate whose unsigned minimal noncommutation count is $E_6$ and whose signed phase
geometry is the protected $81$-dimensional sector.

% ===========================================================
\section{Open problems}
% ===========================================================

\begin{enumerate}[label=(\arabic*)]
\item Derive the rank-$81$ projector without enumeration, directly from the representation
theory of $\operatorname{PSp}(4,3)$ or the Weyl group of $E_6$.
\item Identify the association scheme underlying the $1,3,9,27$ $X$-visibility overlaps.
\item Explain the dual $Z$ distribution $0^{1187},1^{288},2^{96},3^{32},4^{16}$ from a
closed finite-geometric quotient.
\item Construct an explicit equivariant map from nonzero minimal $X/Z$ pairings to the
$51840$ elements of $\Weyl(E_6)$.
\item Identify the transfer operator whose metric generating function is $P(t)$.
\end{enumerate}

\begin{thebibliography}{99}
\bibitem{Goppa1981} V. D. Goppa, \textit{Codes on algebraic curves}, Soviet Math. Dokl. \textbf{24} (1981), 170--172.
\bibitem{Stichtenoth2009} H. Stichtenoth, \textit{Algebraic Function Fields and Codes}, 2nd ed., Springer, 2009.
\bibitem{Brouwer1989} A. E. Brouwer, A. M. Cohen, and A. Neumaier, \textit{Distance-Regular Graphs}, Springer, 1989.
\bibitem{Humphreys1972} J. E. Humphreys, \textit{Introduction to Lie Algebras and Representation Theory}, Springer, 1972.
\bibitem{Adams1996} J. F. Adams, \textit{Lectures on Exceptional Lie Groups}, University of Chicago Press, 1996.
\bibitem{Slansky1981} R. Slansky, \textit{Group theory for unified model building}, Physics Reports \textbf{79} (1981), 1--128.
\bibitem{GeorgiGlashow1974} H. Georgi and S. L. Glashow, \textit{Unity of all elementary-particle forces}, Phys. Rev. Lett. \textbf{32} (1974), 438--441.
\bibitem{Coxeter1973} H. S. M. Coxeter, \textit{Regular Polytopes}, 3rd ed., Dover, 1973.
\bibitem{CalderbankShor1996} A. R. Calderbank and P. W. Shor, \textit{Good quantum error-correcting codes exist}, Phys. Rev. A \textbf{54} (1996), 1098--1105.
\bibitem{Steane1996} A. M. Steane, \textit{Error correcting codes in quantum theory}, Phys. Rev. Lett. \textbf{77} (1996), 793--797.
\end{thebibliography}

\end{document}
